\documentclass[11pt,a4paper]{article}

\usepackage[a4paper,margin=2.5cm]{geometry}
\usepackage[T1]{fontenc}
\usepackage[utf8]{inputenc}
\usepackage{lmodern}
\usepackage{amsmath,amssymb}
\usepackage{booktabs}
\usepackage{graphicx}
\usepackage{setspace}
\usepackage{placeins}
\usepackage{float}
\usepackage{adjustbox}
\usepackage[font=small,labelfont=bf]{caption}
\usepackage[hidelinks]{hyperref}

\onehalfspacing
\graphicspath{{../results/report/figures/}}
\setlength{\parskip}{0.3em}
\renewcommand{\topfraction}{0.9}
\renewcommand{\textfraction}{0.1}
\renewcommand{\floatpagefraction}{0.7}

% All floats are in the appendix and are placed exactly where they are written, so the numbering
% follows the order of the source. The exported tables ask for [t]; this makes them [H] instead.
\let\origtable\table
\let\endorigtable\endtable
\renewenvironment{table}[1][]{\origtable[H]}{\endorigtable}

% Table exported by scripts/11_make_report_assets.py (single line spacing so rows stay compact).
% The exported tables call \resizebox{\textwidth}{!}, which also enlarges narrow tables. Here it
% is replaced by a shrink-only box, so every table keeps its natural \footnotesize text unless it
% is wider than the page.
\newcommand{\tbl}[1]{\begingroup\linespread{1}\selectfont
  \def\resizebox##1##2##3{\adjustbox{max width=##1}{##3}}%
  \input{../results/report/tables/#1.tex}\endgroup}
% Figure exported by the same script: \fig{file-stem}{caption}{width as a fraction of the text}.
\newcommand{\fig}[3]{\begin{figure}[H]\centering
  \includegraphics[width=#3\textwidth,height=0.28\textheight,keepaspectratio]{#1.pdf}\caption{#2}\label{fig:#1}\end{figure}}

\newcommand{\Sig}{\boldsymbol{\Sigma}}
\newcommand{\bone}{\mathbf{1}}

\title{Where Does the Estimation Error Go?\\[0.2em]
\large Markowitz, RMT-denoised, PCA-factor and HRP portfolios on daily S\&P~500 data}
\author{DSE4212/QF4212 Project 2 --- Group [number and names to be filled in]}
\date{15 November 2026}

\begin{document}
\maketitle

%% ---------------------------------------------------------------------------------------------
\section{Overview of the problem}

Mean-variance portfolios built from sample moments perform poorly out of sample: the optimizer puts its
largest positions where the estimation error is largest \cite{michaud1989,markowitz1952}. The project
asks \emph{where the estimation error goes}, and whether methods that clean the covariance matrix give
better out-of-sample portfolios than the sample-covariance plug-in.

In one rolling backtest with identical universe, dates, constraints and costs, we compare the sample
covariance, Ledoit--Wolf shrinkage \cite{lw2004}, RMT eigenvalue clipping \cite{mp1967,laloux1999}, a PCA
factor model and hierarchical risk parity (HRP, \cite{lopezdeprado2016}), with 1/N as the naive benchmark
\cite{demiguel2009}. Two hypotheses were fixed beforehand: \textbf{H1}, denoised covariances give lower
out-of-sample volatility than the sample covariance in the global minimum-variance (GMV) portfolio;
\textbf{H3}, the advantage of denoising is larger without the long-only constraint, because the constraint
already regularizes \cite{jagannathanma2003}. Parameters were fixed a priori and never tuned on the test
period; null results are reported as results. All figures and tables are collected in the Appendix.

\paragraph{Summary of findings.} Equal weight has the best Sharpe ratio, and no Sharpe difference is
statistically detectable. GMV cuts volatility by about a quarter relative to 1/N, but denoising does
\emph{not} improve on sample GMV (PCA and RMT are worse), so H1 is not supported; H3 holds only for
Ledoit--Wolf. In a simulation with a known factor covariance the denoisers do win, by only about 1\% of
variance, and the advantage disappears when the covariance is time-varying. The extensions
(Section~5.7) do not change the picture.

%% ---------------------------------------------------------------------------------------------
\section{Dataset description}

We use only the supplied Kaggle files of daily S\&P~500 prices. The stock file lists 502 symbols, but only
172 have prices (the other 330 are empty; Apple, for example, has none). Cleaning is done in code: symbols
without prices are dropped and logged, and 11 suspect returns across 9 symbols are flagged. The cleaned panel
runs from 2010-01-05 to 2024-12-20 (3{,}767 return dates). The index file is used only for a comparison on its
common period, and sector labels only as labels.

At each rebalance date the universe is the top $N=100$ names by trailing 504-day dollar volume among symbols
with a complete $T=504$-day history, using data up to that date only (150 to 170 symbols are eligible).
The first window ends 2012-01-03, so the out-of-sample period is 2012-01-04 to 2024-12-20 (3{,}263 trading
days, 156 monthly rebalances). The 172 priced names are current index members, so the universe carries
survivorship bias; relative comparisons share it. $N=200$, a planned sensitivity value, is infeasible.

%% ---------------------------------------------------------------------------------------------
\section{Initial data analysis}

Daily returns are strongly fat-tailed: on the first, stress and last snapshot windows the pooled excess
kurtosis is 11.7, 25.9 and 29.2 (Figure~\ref{fig:figA01_return_hist}), so windows are winsorized per column
(0.1\%/99.9\%) before estimation.

For $N$ assets and $T$ observations of pure noise, the eigenvalues of a correlation matrix lie below the
Marchenko--Pastur (MP) edge $\lambda_+=(1+\sqrt{q})^2$, $q=N/T$ \cite{mp1967}; here $q=0.20$ and
$\lambda_+=2.09$. Figure~\ref{fig:fig01_mp_spectrum} shows that only 2, 4 and 6 eigenvalues rise above the
edge at the three snapshots, the largest (the market mode) taking 53\%, 49\% and 23\% of the total; the rest
is a bulk that the sample estimate spreads too widely, which is what RMT clipping targets. Across the 156
windows, the number of signal eigenvalues ranges from 2 to 6 (median 4) and the average correlation from
0.21 to 0.52, highest at the start (Figure~\ref{fig:fig03_spectrum_time}). The 21-day average correlation
is 0.32 over the sample, 0.38 in 2020 and 0.43 in 2022: correlations jump in stress periods, so an estimator
calibrated on a calm window is exposed.

%% ---------------------------------------------------------------------------------------------
\section{Methods used}

The central tools are unsupervised learning methods (eigen-decomposition and PCA, random matrix theory,
hierarchical clustering) that build the covariance input of the optimizer. Supervised learning (Ridge and
random forest) enters in one extension.

\paragraph{Backtest.} Log returns are used for estimation and simple returns for profit and loss. Weights at
rebalance date $t$ use only data up to $t$ (including the universe and winsorization) and apply from $t+1$,
drift between rebalances and pay 10 bps on turnover. Every strategy is long-only, fully invested and faces
the same universe, dates and costs. Unit tests check the absence of look-ahead and the accounting.

\paragraph{Covariance estimators} (written by hand in NumPy/SciPy). \emph{Sample} covariance $S$.
\emph{Ledoit--Wolf}: $\hat\Sig=(1-\delta)S+\delta(\operatorname{tr}S/N)I$ with a data-driven $\delta$
\cite{lw2004}. \emph{RMT clipping}: correlation eigenvalues above $\lambda_+$ are kept, the rest replaced by
their mean (trace preserved), rescaled to a unit diagonal \cite{laloux1999}. \emph{PCA factor model}:
$\hat\Sig=B\Lambda_KB^\top+\operatorname{diag}(\text{residual})$ with $K$ from the MP rule. \emph{HRP}:
clustering of the correlation matrix and recursive bisection by cluster variance, with no matrix inversion
\cite{lopezdeprado2016}.

\paragraph{Portfolios.} 1/N; GMV, $\min_w w^\top\hat\Sig w$ s.t.\ $\bone^\top w=1$, $w\ge0$, for each
covariance; HRP; and long-only maximum-Sharpe (Markowitz) with sample and RMT covariance, using window means
(mean log return plus half the variance) as expected returns and a zero risk-free rate.

\paragraph{Evaluation.} Annualized return, volatility, Sharpe, Sortino, drawdown, CVaR, turnover and the
effective number of names $1/\sum_i w_i^2$, net and gross of costs, plus the volatility the optimizer
\emph{predicts} against the one realized over the next 21 days. Differences are tested with a studentized
circular block bootstrap (block 10, 5{,}000 resamples, strategies resampled jointly; Sharpe differences and log
variance ratios), Holm adjustment within each family and the smallest effect detectable at 80\% power.

\paragraph{Simulation and extensions.} A factor-model simulation with known covariance ($N=100$, three true
factors, 200 replications; stationary and time-varying) separates ``the estimator is worse'' from ``the data
are different''. Extensions use the same engine: inverse-volatility and equal-risk-contribution (ERC)
portfolios \cite{maillard2010}; ridge, shrinkage-to-1/N, L1-bounded and elastic-net weights; quadratic-inverse
nonlinear shrinkage \cite{lw2022}; EWMA covariances; a sector-block residual covariance with a random-label
placebo; and the QLIKE forecast loss of each covariance on random portfolios (Newey--West errors,
\cite{neweywest1987}).

\paragraph{Machine-learning return forecasts.} Ridge and a shallow random forest (200 trees, depth 4) predict the
cross-sectional rank of the next 21-day return from 12-1 month momentum, 1-month return, 63-day volatility and
dollar volume. Training uses an expanding window, a refit every six rebalances and a purge so no training label
reaches past the refit date. The signal is judged by its out-of-sample rank information coefficient (IC); a
pre-set gate (mean IC $>0$ and $t\ge2$) decides whether it may feed the optimizer, as a centered rank score with
the RMT covariance. All hyperparameters were fixed before the run.

%% ---------------------------------------------------------------------------------------------
\section{Model performance evaluation and application results}

\subsection{Estimated versus realized Markowitz frontier}

Plug-in Markowitz treats estimates as truth. At the first window (2012-01-03) the maximum-Sharpe portfolio has an
estimated Sharpe ratio of 2.07 with the sample covariance against 0.87 realized over the next 252 days; at the
latest date with a full year ahead (2023-12-20) the figures are 1.87 and 1.01 (Figure~\ref{fig:fig04_frontier}).
The stress date (2020-03-31) is an exception because of the rebound that followed. The RMT covariance does not
change the picture, which points to expected returns as the dominant error in the Markowitz problem.
\subsection{Out-of-sample performance}

Table~\ref{tab:main_net} and Figure~\ref{fig:fig05_equity} give the main result. Equal weight has the highest
Sharpe ratio (0.99) and HRP is next (0.93 to 0.94). GMV portfolios have the lowest volatility (about 0.13
against 0.17 for 1/N) but also the lowest return, so their Sharpe ratios are lower (0.66 to 0.76). The Markowitz
portfolios earn a return close to 1/N with higher volatility (0.20) and a turnover of about 6 a year, against 0.9 for 1/N and 2.2 to 2.7 for GMV and HRP. Costs of 10 bps lower Sharpe ratios by only
0.01 to 0.03 (Table~\ref{tab:main_gross}). Denoising does not lower GMV volatility: RMT (0.136) and PCA (0.135)
are above sample GMV (0.130), and Ledoit--Wolf matches it. GMV portfolios are concentrated (about 10 effective
names against about 66 for HRP, Table~\ref{tab:concentration}) and tilt towards Consumer Defensive and Utilities
and away from Technology (Figure~\ref{fig:figA11_struct_sectors}).

\subsection{Statistical evidence}

All GMV variants reduce volatility relative to 1/N by 22 to 25\% (significant after Holm). Against sample GMV, PCA
has 3.3\% higher volatility (95\% interval 2.4 to 4.3\%, significant), Ledoit--Wolf is 0.6\% lower (significant but
negligible), and RMT is 4.5\% higher with an interval that includes zero in the studentized test ($p=0.060$). The
RMT result is therefore suggestive, not established: its Holm-adjusted $p$-value is 0.022, 0.060 and 0.201 for block
lengths 5, 10 and 20, while all other conclusions are stable. The Markowitz portfolios have about 1.5 times the
volatility of sample GMV.

\textbf{No Sharpe difference is significant}, but this is mostly a power problem, not evidence of equality: the
smallest detectable difference is about 0.5 to 0.6 against 1/N and 0.13 to 0.17 against sample GMV.
\subsection{Does the optimizer know its own risk?}

The pooled ratio of realized to predicted variance is 1.32 for sample GMV, 1.39 for RMT and 1.61 for PCA, so the
optimizers under-predict their own risk, and the denoised versions by more; HRP does not (0.97)
(Figure~\ref{fig:fig07_vol_forecast}, Table~\ref{tab:risk_forecast}). This is consistent with an optimizer that
chooses weights against its own estimate, an effect a better-conditioned covariance does not remove.

\subsection{Subperiods, sensitivity and the long-only constraint}

Rankings are not stable across subperiods (Figure~\ref{fig:fig09_subperiods}). Minimum-variance portfolios have a
Sharpe ratio near zero in the COVID stress period (0.06 to 0.08 against 0.67 for 1/N) while the Markowitz
portfolios exceed 1.5; in 2022 every strategy is negative, GMV-PCA the least ($-0.02$); in calm markets 1/N and HRP
lead. A 250-day subperiod has a Sharpe standard error near 1, so this is descriptive. Across 15 sensitivity
scenarios (81 strategy-parameter combinations, reported as sensitivity and never used to choose a configuration),
one result is stable: at every $q=N/T$ from 0.10 to 0.40, GMV-RMT and GMV-PCA have higher volatility than sample
GMV (Figure~\ref{fig:fig10_sens_q}). The loss does not go away with more data per asset.

\textbf{H3} (Table~\ref{tab:h3_tests}) holds only for Ledoit--Wolf (difference $-0.055$, interval $-0.080$ to
$-0.033$). For RMT there is no difference, and PCA goes the other way: it is worse than sample GMV, more so
unconstrained. Without the constraint, leverage is 2.7 to 3.7, turnover 6 to 12 times higher and volatility not
lower for any estimator (shorts carry no borrow cost here, which flatters them). The long-only constraint is
already a strong regularizer \cite{jagannathanma2003}, leaving little for denoising to add.

\subsection{What a known-truth simulation says}

On real data the true covariance is unknown, so Figure~\ref{fig:fig12_simulation} repeats the comparison with a
known factor covariance. With a stationary covariance, PCA and RMT beat sample GMV, by the most at small $T$ (log
true-variance ratio about $-0.06$ at $T=126$) but only $-0.012$ (PCA) and $-0.005$ (RMT), about 1\% of variance, at
$T=504$; without the long-only constraint the gain is far larger (about $-1.3$ at $T=126$). With time-varying volatilities and loadings the advantage vanishes for $T\ge504$: the error is
mostly non-stationarity bias, which denoising cannot remove. The simulated loss is only about 0.7\% of variance,
however, far from the 3 to 5\% of volatility (about 7 to 9\% of variance) seen on real data. The engine is not the
cause: run unchanged on synthetic stationary panels, it agrees with the direct evaluation to $10^{-15}$ and
reproduces the simulation's ranking and magnitudes.

\subsection{Why does denoising not help? Extensions}

\paragraph{Non-stationarity.} EWMA covariances (half-lives 126 and 63 days, fixed in advance) win in the simulated
time-varying world; on real data their point estimates beat sample GMV (volatility ratio 0.96, Sharpe 0.83 against
0.74) but not significantly, as the test could only detect a 5 to 10\% change. Non-stationarity is consistent with
the evidence but not confirmed, and it does not explain the RMT/PCA loss (EWMA-RMT is again no better than EWMA).

\paragraph{Sector structure in the residual.} The PCA residual still holds sector correlation: the same-sector minus
cross-sector gap is 0.061 against 0.000 for shuffled labels (permutation $p<0.05$ in all 156 windows, descriptive), so the factors remove only about 45\% of the sector gap. Restoring same-sector blocks lowers PCA's
volatility by 2.4\% (Holm $p=0.001$), removing about three quarters of its loss, but stays above sample GMV (ratio
1.009, Holm $p=0.036$). True sectors beat all 20 random-label draws (rank $p=0.048$, the smallest possible), and
random labels recover only 18\% of the gain. Sector structure is a real but incomplete part of the PCA loss.

\paragraph{Forecast quality versus use.} As forecasters of the variance of random long-only portfolios, RMT and PCA
are not worse than the sample covariance: QLIKE differences are not significant after Holm (smallest adjusted
$p=0.77$), and the point estimates rank PCA-sector, PCA, RMT, sample, nonlinear shrinkage, Ledoit--Wolf identically
in three portfolio families (Table~\ref{tab:forecastloss}). This suggests that the GMV loss comes from how the
optimizer exploits the estimate, not from a worse covariance as such. The mechanism was not tested directly, and
the test has little power.

\paragraph{Other optimizers and regularization.} ERC portfolios are almost insensitive to the covariance estimator
(differences below 1\%) and have realized-to-predicted volatility of about 1.0 against 1.19 to 1.32 for GMV, so the
denoising loss is specific to the minimum-variance optimizer, which exploits the small eigenvalues
(Table~\ref{tab:riskbased_metrics}). Ridge and shrinkage towards 1/N push realized volatility towards the
predicted level (1.19 to 1.01) at the price of higher volatility, with no significant Sharpe gain
(Table~\ref{tab:regularized_metrics}). Nonlinear shrinkage, the theoretical counterpart of RMT clipping, does not
beat sample GMV either (volatility ratio 0.996, Holm $p=0.12$).

\paragraph{Machine-learning return forecasts.} Both models fail the pre-set gate (Table~\ref{tab:ml_signal_ic}): the
mean rank IC is 0.025 for Ridge ($t=1.25$) and 0.029 for the random forest ($t=1.70$) over 131 dates, against a gate
of $t\ge2$, and the single features are no better. The portfolios built on the forecasts are shown for completeness
(Table~\ref{tab:ml_metrics}): the optimizer turns a noise ranking into a concentrated, high-turnover portfolio
(about 5 effective names, turnover 17 to 19 a year) with volatility 15\% to 22\% above MV-RMT (Holm $p=0.02$ and
$0.001$) and a lower Sharpe ratio. This is a null result. The centered-rank mapping drops the market premium that
MV-RMT keeps, so part of the gap may come from the mapping, which was fixed before the results.

%% ---------------------------------------------------------------------------------------------
\section{Conclusion}

\paragraph{Answer to the research question.} On this data the estimation error does \emph{not} go where denoising
can remove it. Cleaning the covariance matrix reduces condition numbers by a factor of three to six, yet does not
lower the out-of-sample volatility of GMV and, for PCA, raises it. Our evidence points to three parts of the
explanation, none isolated fully: sector correlation that denoising throws away (real but incomplete),
non-stationarity of the covariance (consistent with the simulation, underpowered on real data), and the optimizer
exploiting its estimate (covariance forecasts are not worse, but GMV portfolios under-predict their own risk). The
long-only constraint already regularizes much of the problem. In risk-adjusted terms 1/N is hard to beat: GMV gives
the lowest volatility at a lower Sharpe ratio, while ERC and HRP give diversified portfolios with a Sharpe ratio
close to 1/N.

\paragraph{Limitations.} (i) Only 172 of 502 symbols have prices and they are current members, so survivorship bias
remains, the sample is mostly a bull market, and $N=200$ was infeasible. (ii) One historical path: the bootstrap
quantifies sampling noise within it, and many differences are inside the detectable-effect bounds, so ``not
significant'' means ``not detectable''. (iii) Costs are proportional, with no market impact and no borrow cost for
shorting variants. (iv) The ML evaluation starts two years after the core backtest (132 dates), uses four features
and treats the 21-day-spaced IC values as independent. (v) The ``optimizer selects against the estimate'' mechanism
rests on indirect evidence, and the sector and non-stationarity diagnostics are descriptive with overlapping windows.

\paragraph{Future work.} Test the optimizer mechanism directly (for example by resampling diagnostics); repeat the
study on a full point-in-time S\&P~500 file with delistings; combine time-varying volatility with denoising;
denoise eigenvectors as well as eigenvalues; and try richer return signals with a longer training history.

\paragraph{Reproducibility.} Parameters are in \texttt{config.yaml}; the submitted code regenerates every table and figure.

%% ---------------------------------------------------------------------------------------------
\begin{thebibliography}{99}\small
\bibitem{demiguel2009} DeMiguel, V., Garlappi, L. and Uppal, R. (2009). Optimal versus naive diversification:
how inefficient is the 1/N portfolio strategy? \emph{Review of Financial Studies} 22(5), 1915--1953.
\bibitem{jagannathanma2003} Jagannathan, R. and Ma, T. (2003). Risk reduction in large portfolios: why
imposing the wrong constraints helps. \emph{Journal of Finance} 58(4), 1651--1683.
\bibitem{laloux1999} Laloux, L., Cizeau, P., Bouchaud, J.-P. and Potters, M. (1999). Noise dressing of
financial correlation matrices. \emph{Physical Review Letters} 83, 1467--1470.
\bibitem{lw2004} Ledoit, O. and Wolf, M. (2004). A well-conditioned estimator for large-dimensional
covariance matrices. \emph{Journal of Multivariate Analysis} 88(2), 365--411.
\bibitem{lw2022} Ledoit, O. and Wolf, M. (2022). Quadratic shrinkage for large covariance matrices.
\emph{Bernoulli} 28(3), 1519--1547.
\bibitem{lopezdeprado2016} L{\'o}pez de Prado, M. (2016). Building diversification portfolios that outperform
out of sample. \emph{Journal of Portfolio Management} 42(4), 59--69.
\bibitem{maillard2010} Maillard, S., Roncalli, T. and Te{\"i}letche, J. (2010). The properties of equally
weighted risk contribution portfolios. \emph{Journal of Portfolio Management} 36(4), 60--70.
\bibitem{markowitz1952} Markowitz, H. (1952). Portfolio selection. \emph{Journal of Finance} 7(1), 77--91.
\bibitem{mp1967} Marchenko, V. A. and Pastur, L. A. (1967). Distribution of eigenvalues for some sets of
random matrices. \emph{Mathematics of the USSR-Sbornik} 1(4), 457--483.
\bibitem{michaud1989} Michaud, R. O. (1989). The Markowitz optimization enigma: is ``optimized'' optimal?
\emph{Financial Analysts Journal} 45(1), 31--42.
\bibitem{neweywest1987} Newey, W. K. and West, K. D. (1987). A simple, positive semi-definite,
heteroskedasticity and autocorrelation consistent covariance matrix. \emph{Econometrica} 55(3), 703--708.
\end{thebibliography}

%% ---------------------------------------------------------------------------------------------
\clearpage
\appendix
\section{Appendix: figures and tables}

All tables are also available as CSV, Markdown and LaTeX in \texttt{results/report/tables}, and all figures as
PDF and PNG in \texttt{results/report/figures}. The long tables (bootstrap tests, subperiods, sensitivity, the
parameter table) are not repeated here; see \texttt{tab03\_significance}, \texttt{tab04\_subperiods},
\texttt{tab05\_sensitivity}, \texttt{tab08\_parameters} and \texttt{tabA09\_inference}.

\subsection{Data analysis (Section 3)}
\fig{figA01_return_hist}{Pooled return distribution.}{0.7}
\fig{fig01_mp_spectrum}{Eigenvalue spectrum of the sample correlation matrix against the Marchenko--Pastur
density (first, stress and last snapshot windows).}{0.7}
\fig{fig03_spectrum_time}{Number of eigenvalues above the MP edge and the market-mode share over
time.}{0.7}
\FloatBarrier

\subsection{Main results (Sections 5.1 to 5.6)}
\fig{fig04_frontier}{Estimated against realized efficient frontier at three dates (sample and RMT
covariance).}{0.7}
\tbl{tab01_main_net}
\fig{fig05_equity}{Net equity curves of the core strategies, 1/N and the S\&P~500 (common period).}{0.7}
\tbl{tab02_main_gross}
\tbl{tabA02_concentration}
\fig{fig07_vol_forecast}{Realized over predicted volatility at each rebalance, by strategy.}{0.7}
\tbl{tabA03_risk_forecast}
\fig{fig09_subperiods}{Net Sharpe ratio and maximum drawdown by subperiod.}{0.7}
\fig{fig10_sens_q}{Sensitivity to the ratio $q=N/T$ (common window).}{0.7}
\fig{figA08_sens_other}{Sensitivity to costs, weight cap, number of PCA factors and HRP linkage.}{0.7}
\tbl{tab07_h3_tests}
\tbl{tab06_h3_metrics}
\fig{fig11_h3}{Long-only against unconstrained GMV (hypothesis H3).}{0.7}
\fig{fig12_simulation}{Simulation study with Gaussian returns: true variance of each GMV portfolio relative to the
oracle, by window length.}{0.7}
\fig{figA13_simulation_t}{Simulation study with Student-$t$ returns.}{0.7}
\tbl{tabA01_common_period}
\fig{figA11_struct_sectors}{Mean sector exposure of each strategy (current sector labels).}{0.7}
\FloatBarrier

\subsection{Extensions (Section 5.7)}
\tbl{tabA11_nonstat_metrics}
\tbl{tabA13_nonstat_calibration}
\tbl{tabA17_sectorblock_metrics}
\fig{figA14_res_blocks}{Mean PCA-residual correlation by sector pair.}{0.7}
\fig{figA15_res_timeseries}{Sector structure of the PCA residual over time.}{0.7}
\tbl{tabA20_forecastloss}
\tbl{tabA21_nls_metrics}
\tbl{tabA24_riskbased_metrics}
\tbl{tabA26_regularized_metrics}
\tbl{tabA28_ml_signal_ic}
\tbl{tabA29_ml_metrics}

\end{document}
